%% T6.tex -- Proposition 6: Mixing bound / approach rate to R_infinity.
%%
%% REVISED (2026-04-22):
%%   * Replaced the informal "generically diagonalizable" handwave with
%%     (a) an EXPLICIT, norm-agnostic bound for the diagonalizable case
%%     using the eigenvector condition number kappa_infty(V), and (b) a
%%     formal Jordan-canonical-form bound for the non-diagonalizable
%%     case, stating the polynomial factor pi(d;k) arising from the
%%     nilpotent super-diagonal of a size-k Jordan block.
%%   * The practical horizon formula (Eq. t6-steps) is now stated in
%%     two forms: the clean closed form for k=1, and an implicit
%%     inequality for general k with an explicit upper bound on the
%%     required horizon inflation.
%%   * Citations: Horn-Johnson for Jordan-form norm bounds; Stewart
%%     for the condition-number-of-eigenvectors bound.

\noindent\emph{The following is a direct application of classical
norm bounds on powers of a matrix with spectral radius below one
\cite[Ch.~3]{hornjohnson2013} and of the eigenvector-conditioning
bound \cite[Ch.~1]{stewart1998}, specialized to the absorbing
agent-chain setting. The mathematical content is textbook; the
contribution is the \emph{practical horizon} formula derived below,
which gives an actionable ``how many steps is enough?'' answer for
agent deployments.}

\paragraph{Norm and matrix conventions.}
All matrix norms below are the induced $\infty$-norm
$\|A\|_\infty = \max_i \sum_j |A_{ij}|$ unless otherwise noted.
For a diagonalizable matrix $Q = V\Lambda V^{-1}$ we write
$\kappa_\infty(V)\triangleq \|V\|_\infty\|V^{-1}\|_\infty$ for the
$\infty$-norm condition number of the eigenvector matrix; for a
general matrix with Jordan decomposition $Q = P J P^{-1}$ we write
$\kappa_\infty(P) \triangleq \|P\|_\infty\|P^{-1}\|_\infty$.  Let $k$
denote the size of the largest Jordan block associated with an
eigenvalue of magnitude $\rho(Q)$; $k=1$ exactly when $Q$ is
diagonalizable.  Define the Jordan polynomial
\begin{equation}
\pi(d;k) \;\triangleq\; \sum_{j=0}^{\min(k-1,d)} \binom{d}{j},
\qquad d\ge 0,
\label{eq:jordan-poly}
\end{equation}
which satisfies $\pi(d;1)=1$ and, for $d\ge k-1$,
$\pi(d;k)\le k\,\binom{d}{k-1}\le \tfrac{k}{(k-1)!}(d+1)^{k-1}$.

\begin{proposition}[Approach-Rate Bound; adapted from Horn--Johnson
  \cite{hornjohnson2013} and Stewart \cite{stewart1998}]
\label{thm:mixing}
Under Assumptions~\ref{A:trans}--\ref{A:init} with $\rho\triangleq
\rho(Q)<1$ and $N=(I-Q)^{-1}$, the absorbing-success gap admits the
following explicit, norm-agnostic bounds.
\begin{enumerate}[(i)]
\item \emph{(Diagonalizable case, $k=1$.)} If $Q = V\Lambda V^{-1}$,
then
\begin{equation}
R_\infty - R(d)
  \;\le\; \kappa_\infty(V)\,\rho^{d}\,\bigl\|N R_\oplus\bigr\|_\infty,
  \qquad d\ge 0,
\label{eq:t6-diag}
\end{equation}
and the minimum horizon to bring the gap below $\delta>0$ is
\begin{equation}
  d^\star_{\mathrm{diag}}(\delta)
  \;=\;
  \Bigl\lceil
   \frac{\log\bigl(\kappa_\infty(V)\,\|NR_\oplus\|_\infty / \delta\bigr)}{-\log\rho}
  \Bigr\rceil.
\label{eq:t6-steps-diag}
\end{equation}

\item \emph{(General case via Jordan form.)} For an arbitrary $Q$
with Jordan decomposition $Q=PJP^{-1}$ and largest leading
Jordan-block size $k$,
\begin{equation}
R_\infty - R(d)
  \;\le\;
  \kappa_\infty(P)\,\pi(d;k)\,\rho^{d-k+1}\,\bigl\|N R_\oplus\bigr\|_\infty,
  \qquad d\ge k-1.
\label{eq:t6-jordan}
\end{equation}
The horizon to reach gap $\delta$ is then the smallest $d\ge k-1$
satisfying
\begin{equation}
\kappa_\infty(P)\,\pi(d;k)\,\rho^{d-k+1}\,\|NR_\oplus\|_\infty \;\le\; \delta,
\label{eq:t6-steps-jordan}
\end{equation}
which admits the \emph{explicit upper bound}
\begin{equation}
  d^\star(\delta,k)
  \;\le\;
  \Bigl\lceil
    \tfrac{1}{-\log\rho}\,\log\!\tfrac{\kappa_\infty(P)\,k\,\|NR_\oplus\|_\infty}{\delta}
    + (k-1)\,\bigl(\tfrac{\log(d_0+1)}{-\log\rho}+1\bigr)
  \Bigr\rceil,
\label{eq:t6-jordan-bound}
\end{equation}
where $d_0$ is the diagonalizable horizon \eqref{eq:t6-steps-diag}
(with $\kappa_\infty(P)$ in place of $\kappa_\infty(V)$). Thus the
polynomial Jordan factor inflates the required horizon by an additive
$(k{-}1)\log(d_0{+}1)/(-\log\rho) + (k{-}1)$ steps, not multiplicatively.
\end{enumerate}
\end{proposition}

%% ----------------------------------------------------------------
%% [REDUCTION 2026-04-24] Proof and remarks suppressed from body.
%% ----------------------------------------------------------------
\iffalse
\begin{proof}
\emph{Step 1 --- gap formula (unchanged).}
From Def.~\ref{def:reliabinf} and Thm.~\ref{thm:closed},
$R_\infty - R(d) = e_{s_0}^\top Q^d N R_\oplus$. Hence
\[
\bigl|R_\infty - R(d)\bigr| \;\le\; \|Q^d\|_\infty\,\|NR_\oplus\|_\infty,
\]
where we used $\|e_{s_0}^\top\|_\infty=1$. The task is therefore an
explicit bound on $\|Q^d\|_\infty$.

\emph{Step 2 --- diagonalizable bound (i).}
If $Q=V\Lambda V^{-1}$, then $Q^d = V\Lambda^d V^{-1}$ and
sub-multiplicativity of an induced norm gives
\[
\|Q^d\|_\infty \;\le\; \|V\|_\infty\,\|\Lambda^d\|_\infty\,\|V^{-1}\|_\infty
                \;=\; \kappa_\infty(V)\,\|\Lambda^d\|_\infty.
\]
The matrix $\Lambda^d=\operatorname{diag}(\lambda_1^d,\dots,\lambda_m^d)$ is
diagonal, so $\|\Lambda^d\|_\infty = \max_i|\lambda_i|^d = \rho^d$.
Substituting yields \eqref{eq:t6-diag}. Solving the resulting
inequality for $d$ yields the closed form \eqref{eq:t6-steps-diag}.
\emph{This step no longer uses ``generic diagonalizability'' as a
handwave}: the $\kappa_\infty(V)$ factor is explicit and finite for
every diagonalizable $Q$, and is $1$ exactly when $V$ is orthogonal
(the classical normal-matrix case).

\emph{Step 3 --- Jordan-form bound (ii).}
For general $Q$, let $Q = P J P^{-1}$ with $J$ a direct sum of Jordan
blocks $J_{k_i}(\lambda_i)$. The block
$J_k(\lambda) = \lambda I_k + N_k$, where $N_k$ is the $k\times k$
nilpotent super-diagonal shift, has $N_k^k = 0$, so the binomial
expansion truncates:
\begin{equation}
J_k(\lambda)^d
  \;=\; \sum_{j=0}^{\min(k-1,d)} \binom{d}{j}\lambda^{d-j}\,N_k^{\,j}.
\label{eq:jordan-expansion}
\end{equation}
Each $N_k^{\,j}$ has entries in $\{0,1\}$ only, with a single $1$ on
the $j$-th super-diagonal and zeros elsewhere; hence $\|N_k^{\,j}\|_\infty
= 1$ for $0\le j\le k-1$. Taking $\infty$-norms and summing,
\begin{align*}
\bigl\|J_k(\lambda)^d\bigr\|_\infty
 &\;\le\; \sum_{j=0}^{\min(k-1,d)} \binom{d}{j}\,|\lambda|^{d-j} \\
 &\;\le\; |\lambda|^{d-k+1}\,\sum_{j=0}^{k-1}\binom{d}{j}\,|\lambda|^{k-1-j}
  \;\le\; \rho^{d-k+1}\,\pi(d;k),
\end{align*}
where the last inequality uses $|\lambda|\le\rho\le 1$ in the second
factor (so $|\lambda|^{k-1-j}\le 1$) and the definition of $\pi(d;k)$ in
\eqref{eq:jordan-poly}. Taking $\infty$-norms at the block-diagonal
level and using sub-multiplicativity through $\kappa_\infty(P)$
finally gives
\[
\|Q^d\|_\infty \;\le\; \kappa_\infty(P)\,\pi(d;k)\,\rho^{d-k+1},
\qquad d\ge k-1,
\]
which combined with Step 1 yields \eqref{eq:t6-jordan}.

\emph{Step 4 --- horizon inflation.}
Starting from \eqref{eq:t6-steps-jordan}, take logarithms:
\[
\log\kappa_\infty(P) + \log\pi(d;k) + (d-k+1)\log\rho
  \;\le\; \log(\delta/\|NR_\oplus\|_\infty).
\]
Using $\pi(d;k)\le k\,\binom{d}{k-1} \le k(d+1)^{k-1}/(k-1)!$ and
the monotonicity of $\log(1+x)/x$, the rightmost inequality gives
$\log\pi(d;k) \le \log k + (k-1)\log(d+1)$. Rearranging and using
$-\log\rho>0$:
\[
d \;\le\; k-1 + \frac{1}{-\log\rho}\,
            \Bigl[\log\tfrac{\kappa_\infty(P)\,k\,\|NR_\oplus\|_\infty}{\delta}
               + (k-1)\log(d+1)\Bigr].
\]
Substituting any fixed $d\ge d_0$ (the diagonalizable horizon) on
the right-hand side gives an upper bound on the smallest $d$ that
satisfies the inequality, which is \eqref{eq:t6-jordan-bound}. The
additive form confirms that the Jordan block inflates the required
horizon by at most an additive term of order $(k-1)\log d/(-\log\rho)$,
\emph{not} multiplicatively --- a small correction for the
$\rho\in[0.5,0.9]$ regime typical of agent chains, in which a size-$2$
Jordan block adds roughly $2$--$5$ steps.
\end{proof}

\begin{remark}[When is $k=1$ guaranteed?]
$k=1$ exactly when $Q$ is diagonalizable. This holds whenever $Q$
has $m$ distinct eigenvalues (generic perturbation argument) but
\emph{is not} automatic for structured agent chains. In particular,
forward-chain agents with a deterministic sub-goal ladder (fixed
\texttt{while} loops, sequential API calls, or scripted
branches) typically produce a strictly upper-triangular nilpotent
component in $Q$, for which eigenvalues of magnitude $\rho$ can
\emph{coalesce} and produce Jordan blocks of size $k\ge 2$.
Proposition~\ref{thm:mixing}(ii) is the bound that applies to those
chains; (i) is only a shortcut when the practitioner has verified
diagonalizability (which can be tested in $O(m^3)$ by the numerical
eigendecomposition routine already used in
\S\ref{sec:sim}).
\end{remark}

\begin{remark}[Normal chains; the classical case]
When $Q$ is symmetric (or more generally normal) the eigenvector
matrix $V$ may be chosen orthogonal, in which case
$\kappa_\infty(V)\le\sqrt{m}$ by norm-equivalence (orthogonal
matrices have $\|V\|_2 = 1$ and $\|V\|_\infty \le \sqrt{m}\|V\|_2$)
and \eqref{eq:t6-diag} reduces to the classical spectral bound
$\|Q^d\|_2 \le \rho(Q)^d$, recovering the folklore statement that
the original draft implicitly invoked. The new form
\eqref{eq:t6-diag} makes the $\kappa_\infty(V)$ cost of a
non-normal diagonalization explicit.
\end{remark}

\begin{remark}[Practitioner rule of thumb]
For the diagonalizable case, \eqref{eq:t6-steps-diag} gives the
minimum-horizon recommendation directly; the only new cost over the
original draft is computing $\kappa_\infty(V)$ (one eigendecomposition
of $\widehat{Q}$). For agent chains with $\rho(Q)\in[0.5,0.9]$ and
$\kappa_\infty(V)\le 10$, the horizon is between $7$ and $55$ steps for
$\delta=0.01$. For the Jordan case with $k=2$, the inflation is
typically $2$--$5$ additional steps (from the additive polynomial
term); $k\ge 3$ is empirically rare in agent chains we examined and
can be detected from the eigendecomposition residual
$\|Q - V\Lambda V^{-1}\|$.
\end{remark}

\begin{remark}[Connection to mixing times]
The bound \eqref{eq:t6-jordan} is the absorbing-chain analogue of
the classical ergodic mixing bound
$\|\mu_t-\pi\|_\mathrm{TV}\le C_0\,\rho^t$, with
``mixing'' replaced by ``convergence of the cumulative success
probability to its absorbing-state limit''. The explicit $\kappa(V)$
and Jordan-polynomial factors in \eqref{eq:t6-diag} and
\eqref{eq:t6-jordan} parallel the ``pre-mixing constant'' appearing
in the non-reversible Markov chain literature
\cite[Thm.~12.3]{levin2017markov}.
\end{remark}
\fi
